% page 255 du fac-simile (image p-254 du scan)
% bloc 162.6 x 242.0 mm ; marge gauche 39.4 mm
\pgc{17.780mm}{0.000mm}{
\l{\cel{52}247}
\l{}
\l{\sou{jar}o. Granda vazo de terakotajo, kun du ansi, en qua onu}
\l{\cel{2}konservas aquo, oleo. - EFIS.}
\l{}
\l{jardiniero.Orno-moblo qua suportas kesto provizita ye flori.}
\l{\cel{2}- II.Disho ek diversa sorti di legumi. - DEFIS.}
\l{}
\l{jargonar.(trans.)\cel{3}I.Dicar ulo ye maniero nekomprenebla pro}
\l{\cel{2}nekorekteso vortala, gramatikala, pronuncala, e c. - II.}
\l{\cel{2}Parolar idiomo stranjera koram ulu qua judikas olu kom bar-}
\l{\cel{2}bara pro ne komprenar olu. - DEFIRS.}
\l{}
\l{\sou{jasmino}. (\sou{bot.)}\cel{2}Arbusto sarmentoza, genero tipo di la fa-}
\l{\cel{2}milio \textquotedbl{}jasminei\textquotedbl{}, di qua la flori esas odoroza. - L.}
\l{\cel{2}\sou{jasminus}. - DEFIRS.}
\l{}
\l{\sou{jasp}o.Lapido harda ed opaka, sama-natura kam agato, koloroza}
\l{\cel{2}ye reda, flava o verda, uniforme, bendope o makulope. - DEFIS.}
\l{}
\l{\sou{javelin}o.Armo dardo, longa e dina,uzita da la antiqui, qui}
\l{\cel{2}lansis olu manue. - EFIS.}
\l{}
\l{jeo.(\sou{zool}.)Ucelo, genero \textquotedbl{}paseri konobeka\textquotedbl{}, vicina ye la korvi.}
\l{\cel{2}- L. oarroloa ElandariuE. - EF.}
\l{}
\l{\sou{jele}o.(koquarto). I.Suko de substanco animalala,qua solideskis}
\l{\cel{2}per koldesko. - II. Suko konjelita de frukti. - DEFIS.}
\l{}
\l{\sou{jelozi}o.Shutro, ma ek planketi paralela, horizontala, qui}
\l{\cel{2}povas rotacar cirkum sua axo longesala, quan onu muntas}
\l{\cel{2}avan fenestro, quan onu povas levar o deslevar per kordo}
\l{\cel{2}o kordono ; kande olu esas tote levita, la lati esas apli-}
\l{\cel{2}kita ici sur iti ; e dop qua onu povas vidar sen videsar. -}
\l{\cel{2}DEFIS.}
\l{}
\l{\sou{jema}r.(\sou{netrans.}) Agar lamento ne-artikulita. - FIS.}
\l{}
\l{\sou{jemel}o.I. Qua naskis sama-parture kam altra filio (o yuno). -}
\l{\cel{2}II. (\sou{plurale}). Kastor e Polux,un ek la 22 signi di zodiako.-}
\l{\cel{2}EFIS.}
\l{}
\l{\sou{jena}r. (\sou{trans.}) Igar nekomfortoza, nekomoda, per presar, kle-}
\l{\cel{2}mar, opresar. - DF.}
\l{}
\l{\sou{jendarm}o.Soldato qua apartenas a korpo di qua la tasko kon-}
\l{\cel{2}sistas ye sekurigar la vivo di la civitani. - DEFIS.}
\l{}
\l{\sou{jenio.(armeo)}.\cel{2}La korpo ek la oficiri e soldati qui aplikas}
\l{\cel{2}la cienci a la fortifiko di la fortresi, a la konstrukto}
\l{\cel{2}di navi, di ponti, e c.\cel{2}- DFI.}
\l{}
\l{\sou{jenjiv}o.Tisuo karna qua garnisas la arkadi dentala e la}
\l{\cel{2}perimetro di la denti per adherar a lia kolo. - FI}
\l{}
\l{\sou{jenjivit}o. (\sou{patol}.) Inflameso di la jenjivi. - EFI}
}
